\documentclass{article}
\usepackage{amsmath,amssymb}
\newtheorem{theorem}{Theorem}[section]
\newtheorem{lemma}[theorem]{Lemma}
\newtheorem{definition}{Definition}
\begin{document}
\maketitle
\begin{abstract}
Abstract text with $x+y$.
\end{abstract}
\section{Intro}
Text with \(a_b\) inline and a citation \cite{Ref1} plus \ref{eq:num} and \eqref{eq:num2}.
\begin{equation}\label{eq:num}
  x = \frac{1}{2}
\end{equation}
\begin{align}\label{eq:num2}
  a &= b \
  c &= d
\end{align}
\[
  y = z \tag{star}
\]
\begin{equation*}
  p = q
\end{equation*}
\begin{theorem}[Main]\label{thm:m}
  Something holds for all $n$.
\end{theorem}
\begin{lemma}\label{lem:l}
  Aux.
\end{lemma}
\begin{proof}
  Trivial since \ref{thm:m}.
\end{proof}
\begin{figure}
  \centering
  \begin{tikzpicture}
    \node (a) at (0,0) {A};
    \draw (a) -- (1,1) node[right] {B};
  \end{tikzpicture}
  \caption{A caption. \label{fig:f}}
\end{figure}
\begin{table}
  \caption{T. \label{tab:t}}
  \begin{tabular}{lcr}
    \hline
    A & B \ \hline
    1 & 2 \
    \hline
  \end{tabular}
\end{table}
\begin{itemize}
  \item one
  \item two
\end{itemize}
Text\footnote{A note.} more.

\subsection{Sub}
\begin{enumerate}
 \item first
\end{enumerate}

ewpage
Epsilon \[2mm] x

\appendix
\section{Appendix A}
Bla.
\end{document}
